% Freie Variablen werden aus der vorhandenen Baumrekursion gewonnen.
\subsection{Freie Variablen und Koinzidenz}

In diesem Abschnitt sind \(M=(D,E)\) eine Mengenstruktur und
\(s,t\in A_D\). Variablen sind ihre Codes in \(\mathbb N\).
Alle Aussagen ueber Erfuellung gehoeren zur untersuchenden Theorie.
Insbesondere wird keine Eigenschaft der syntaktischen Beweisbarkeit
vorausgesetzt.

\FormulaDefDeltaK[Operatoren fuer freie Variablen]
{\begin{aligned}
 I(p)&\coloneq\{i\in\mathbb N\mid p=\mathsf q_i\},\\
 L((k,(i,j)))&\coloneq
   \BXLVIIIIf{k=t_{\rm e}\lor k=t_{\rm m}}{\{i,j\}}{\varnothing},\\
 Y&\coloneq\mathcal T\times\powerset(\mathbb N),\\
 f(c)&\coloneq(\mathsf l(c),L(c)),\\
 g((p,U),(q,V))&\coloneq
   (\mathsf n(p,q),U\cup(V\setminus I(p))).
\end{aligned}}
{B48FVOperators}{
 \DeltaRow{Bereiche}{c\in\Sigma,\quad p,q\in\mathcal T}
 \DeltaRow{Wertemengen}{U,V\subseteq\mathbb N}
}
Hier und bei den folgenden rekursiven Konstruktionen bezeichnet eine
Funktionsgleichung \(f(c)\coloneq u(c)\) den beschraenkten Graphen
\(\{(c,z)\in\Sigma\times Y\mid z=u(c)\}\).
Seine Funktionseigenschaft folgt aus
\FormulaRefAuto{TypedTermGraphFunction}[thm], sobald die Werte typisiert
sind; seine Werte liefert \FormulaRefAuto{TypedTermGraphValues}[thm].
Die Hilfsbaeume \(\mathsf n_0,\mathsf c_0,\mathsf q_i\) tragen keine
freien Variablen.

\FormulaThmDeltaK[Konstruktion der Freivariablenfunktion]
{\begin{gathered}
 \exists!V\colon\mathcal T\to\powerset(\mathbb N)\quad
 V(\mathsf l(c))=L(c),\\
 V(\mathsf n(p,q))=V(p)\cup(V(q)\setminus I(p))
\end{gathered}}
{B48FVExists}{
 \DeltaRow{Universell in den Gleichungen}{c\in\Sigma,\quad p,q\in\mathcal T}
}
\begin{tabproofsplitwide}
 \proofpartwide{Typisierung und Rekursion}
 \proofstepwidestar[]{L(c)\subseteq\mathbb N}
  {\FormulaRefAuto{B48FVOperators}[def]}
 \proofstepwidestar[]{U\cup(V\setminus I(p))\subseteq\mathbb N}
  {\FormulaRefAuto{B48FVOperators}[def]}
 \proofstepwidestar[]{f\colon\Sigma\to Y}
  {\FormulaRefAuto{TypedTermGraphFunction}[thm]{1},
   \FormulaRefAuto{B48LeafType}[thm]}
 \proofstepwidestar[]{g\colon Y\times Y\to Y}
  {\FormulaRefAuto{TypedTermGraphFunction}[thm]{2},
   \FormulaRefAuto{B48NodeType}[thm]}
 \proofstepwidestar[]{\begin{gathered}
  \exists!e\colon\mathcal T\to Y\quad
  e(\mathsf l(c))=f(c),\\
  e(\mathsf n(p,q))=g(e(p),e(q))
 \end{gathered}}
  {\FormulaRefAuto{FullPlanarBinaryTreeRecursion}[thm]{3,4}}
 \closeproofpartwide
\end{tabproofsplitwide}
Fixiere die eindeutig bestimmte Auswertung \(e\).
Die erste Komponente ist der ausgewertete Baum selbst. Dies wird
ausdruecklich durch Strukturinduktion gezeigt. Im folgenden Teilbeweis
steht \(H(p)\) fuer \(\exists U\subseteq\mathbb N\;(e(p)=(p,U))\).
\begin{tabproofsplitwide}
 \proofpartwide{Erhaltung des Baumcodes}
 \proofstepwidestar[]{e(\mathsf l(c))=(\mathsf l(c),L(c))}
  {\FormulaRefAuto{B48FVOperators}[def],
   \FormulaRefAuto{FullPlanarBinaryTreeRecursion}[thm]}
 \proofstepwidestar[]{H(\mathsf l(c))}{\rEI{1}}
 \proofstepwidestar[3]{H(p)\land H(q)}{\rA}
 \proofstepwidestar[4]{e(p)=(p,U)\land e(q)=(q,V)}{\rA}
 \proofstepwidestar[4]{\begin{aligned}
 e(\mathsf n(p,q))={}&
 (\mathsf n(p,q),U\cup(V\setminus I(p)))
 \end{aligned}}
  {\FormulaRefAuto{B48FVOperators}[def]{4},
   \FormulaRefAuto{FullPlanarBinaryTreeRecursion}[thm]}
 \proofstepwidestar[4]{H(\mathsf n(p,q))}{\rEI{5}}
 \proofstepwidestar[3]{H(\mathsf n(p,q))}{\rEEStar{3,4:6}}
 \proofstepwidestar[]{H(p)\land H(q)\to H(\mathsf n(p,q))}
  {\rRI{3,7}}
 \proofstepwidestar[]{\forall p\in\mathcal T\;H(p)}
  {\FormulaRefAuto{FullPlanarBinaryTreeStructuralInduction}[thm]
   {\rUI{2},\rUI{8}}}
 \closeproofpartwide
\end{tabproofsplitwide}
In den Zeugenzeilen sind \(U,V\subseteq\mathbb N\) die durch \(H\)
beschraenkten Zeugen. Zwei solche Werte sind gleich, denn aus
\((p,U)=(p,V)\) folgt \(U=V\) durch
\FormulaRefAuto{(a,b)=(c,d)\eqvdash a=c\land b=d}[thm].
Daher ist
\[
 V=\{(p,U)\in\mathcal T\times\powerset(\mathbb N)\mid e(p)=(p,U)\}
\]
nach \FormulaRefAuto{UniqueValuedGraphFunction}[thm] eine Funktion.
Einsetzen von \(e(p)=(p,V(p))\) und \(e(q)=(q,V(q))\) in die beiden
Rekursionsgleichungen liefert die behaupteten Gleichungen.
Eindeutigkeit folgt ebenfalls durch Bauminduktion: Zwei Kandidaten
stimmen an jedem Blatt durch \(L(c)\), und am Knoten nach Einsetzen
der beiden Induktionsgleichheiten ueberein. Funktionsextensionalitaet
\FormulaRefAuto{FunctionExtensionality}[thm] liefert ihre Gleichheit.
Damit ist die folgende Kennzeichnung gerechtfertigt.

\FormulaDefDeltaK[Freie Variablen eines Formelcodes]
{\operatorname{FV}(p)\coloneq V(p)}
{B48FreeVariables}{
 \DeltaRow{Funktion}{V\text{ aus }\FormulaRefAuto{B48FVExists}[thm]}
 \DeltaRow{Baumcode}{p\in\mathcal T}
}
Die Einschraenkung auf \(\mathcal F\) ist die Freivariablenfunktion
der Objektformeln. Die Auswertung auf Hilfsbaeumen bleibt lediglich
ein Hilfsmittel der Konstruktion.

\FormulaThmDeltaK[Rekursionsklauseln fuer freie Variablen]
{\begin{aligned}
 \operatorname{FV}(\mathsf e(i,j))&=\{i,j\},&
 \operatorname{FV}(\mathsf m(i,j))&=\{i,j\},\\
 \operatorname{FV}(\mathsf b)&=\varnothing,&
 \operatorname{FV}(\mathsf{neg}(p))&=\operatorname{FV}(p),\\
 \operatorname{FV}(\mathsf{and}(p,q))
  &=\operatorname{FV}(p)\cup\operatorname{FV}(q),\\
 \operatorname{FV}(\mathsf{ex}_i(p))
  &=\operatorname{FV}(p)\setminus\{i\}.
\end{aligned}}
{B48FVClauses}{
 \DeltaRow{Formelcodes}{p,q\in\mathcal F}
 \DeltaRow{Indizes}{i,j\in\mathbb N}
}
\begin{tabproof}
 \proofstep{}{\operatorname{FV}(\mathsf e(i,j))=\{i,j\}}
  {\FormulaRefAuto{B48FVExists}[thm],
   \FormulaRefAuto{B48FVOperators}[def]}
 \proofstep{}{\operatorname{FV}(\mathsf m(i,j))=\{i,j\}}
  {\FormulaRefAuto{B48FVExists}[thm],
   \FormulaRefAuto{B48FVOperators}[def]}
 \proofstep{}{\operatorname{FV}(\mathsf b)=\varnothing}
  {\FormulaRefAuto{B48FVExists}[thm],
   \FormulaRefAuto{B48TagNe05}[thm],
   \FormulaRefAuto{B48TagNe15}[thm]}
 \proofstep{}{\begin{gathered}
  \operatorname{FV}(\mathsf n_0)=
  \operatorname{FV}(\mathsf c_0)=
  \operatorname{FV}(\mathsf q_i)=\varnothing
 \end{gathered}}
  {\FormulaRefAuto{B48FVOperators}[def],
   \FormulaRefAuto{B48FVExists}[thm]}
 \proofstep{}{I(\mathsf n_0)=I(\mathsf c_0)=\varnothing}
  {\FormulaRefAuto{B48QNeNeg}[thm],
   \FormulaRefAuto{B48QNeAnd}[thm]}
 \proofstep{}{I(\mathsf n(\mathsf c_0,p))=\varnothing}
  {\FormulaRefAuto{B48LeafNotNode}[thm]}
 \proofstep{}{I(\mathsf q_i)=\{i\}}
  {\FormulaRefAuto{B48QuantifierRootInjective}[thm],\rII}
 \proofstep{}{\operatorname{FV}(\mathsf{neg}(p))=\operatorname{FV}(p)}
  {\FormulaRefAuto{B48FVExists}[thm]{4,5}}
 \proofstep{}{\operatorname{FV}(\mathsf n(\mathsf c_0,p))
                  =\operatorname{FV}(p)}
  {\FormulaRefAuto{B48FVExists}[thm]{4,5}}
 \proofstep{}{\operatorname{FV}(\mathsf{and}(p,q))
                  =\operatorname{FV}(p)\cup\operatorname{FV}(q)}
  {\FormulaRefAuto{B48FVExists}[thm]{9,6}}
 \proofstep{}{\operatorname{FV}(\mathsf{ex}_i(p))
                  =\operatorname{FV}(p)\setminus\{i\}}
  {\FormulaRefAuto{B48FVExists}[thm]{4,7}}
\end{tabproof}


\FormulaThmDeltaK[Gleichheit der Quantorenmarken]
{\mathsf q_i=\mathsf q_j\leftrightarrow i=j}
{B48QuantifierRootEquality}{\DeltaRow{Indizes}{i,j\in\mathbb N}}
\begin{tabproof}
 \proofstep{1}{\mathsf q_i=\mathsf q_j}{\rA}
 \proofstep{1}{i=j}{\FormulaRefAuto{B48QuantifierRootInjective}[thm]{1}}
 \proofstep{}{(\mathsf q_i=\mathsf q_j)\to i=j}{\rRI{1,2}}
 \proofstep{4}{i=j}{\rA}
 \proofstep{}{\mathsf q_i=\mathsf q_i}{\rII}
 \proofstep{4}{\mathsf q_i=\mathsf q_j}{\rIE{4,5}}
 \proofstep{}{i=j\to\mathsf q_i=\mathsf q_j}{\rRI{4,6}}
 \proofstep{}{\mathsf q_i=\mathsf q_j\leftrightarrow i=j}{\rLRI{3,7}}
\end{tabproof}

\FormulaDefDeltaK[Saetze der Mengensprache]
{\operatorname{Sent}(\mathcal L_\in)\coloneq
 \{p\in\mathcal F\mid\operatorname{FV}(p)=\varnothing\}}
{B48SentenceCodes}{\DeltaRow{Formelmenge}{\mathcal F}}

\FormulaThmDeltaK[Mitgliedschaft nach Entfernen eines anderen Elements]
{j\in U,\quad j\neq i\ \vdash\ j\in U\setminus\{i\}}
{B48MemberDifferenceSingleton}{\DeltaRow{Mengen}{U,i,j}}
\begin{tabproof}
 \proofstep{1}{j\in U}{\rA}
 \proofstep{2}{j\neq i}{\rA}
 \proofstep{}{j\in\{i\}\leftrightarrow j=i}
  {\FormulaRefAuto{x\in\{a\}\eqvdash x=a}[thm]}
 \proofstep{2}{j\notin\{i\}}{\rLRS{3,2}}
 \proofstep{1,2}{j\in U\land j\notin\{i\}}{\rAI{1,4}}
 \proofstep{1,2}{j\in U\setminus\{i\}}
  {\FormulaRefAuto{x\in A\setminus B\eqvdash x\in A\land x\notin B}[thm]{5}}
\end{tabproof}

\FormulaThmDeltaK[Fehlen ausserhalb des entfernten Elements]
{x\notin U\setminus\{i\},\quad x\neq i\ \vdash\ x\notin U}
{B48MissingFromDifference}{\DeltaRow{Mengen}{U,i,x}}
\begin{tabproof}
 \proofstep{1}{x\notin U\setminus\{i\}}{\rA}
 \proofstep{2}{x\neq i}{\rA}
 \proofstep{3}{x\in U}{\rA}
 \proofstep{2,3}{x\in U\setminus\{i\}}
  {\FormulaRefAuto{B48MemberDifferenceSingleton}[thm]{3,2}}
 \proofstep{1,2,3}{\bot}{\rBI{4,1}}
 \proofstep{1,2}{x\notin U}{\rCI{3,5}}
\end{tabproof}

\FormulaDefDeltaK[Uebereinstimmung auf einem Variablenbereich]
{s\equiv_U t\coloneq\forall j\in U\;(s(j)=t(j))}
{B48Agreement}{
 \DeltaRow{Belegungen}{s,t\in A_D}
 \DeltaRow{Variablenbereich}{U\subseteq\mathbb N}
}

\FormulaThmDeltaK[Punktweise Werte eines Belegungswechsels]
{\begin{aligned}
 s[i\mapsto a](i)&=a,\\
 j\neq i&\ \vdash\ s[i\mapsto a](j)=s(j)
\end{aligned}}
{B48UpdateValues}{
 \DeltaRow{Belegung}{s\in A_D}
 \DeltaRow{Indizes und Wert}{i,j\in\mathbb N,\quad a\in D}
}
\begin{tabproof}
 \proofstep{}{\forall k\in\mathbb N\;
   \BXLVIIIIf{k=i}{a}{s(k)}\in D}
  {\FormulaRefAuto{B48IfTyped}[thm],
   \FormulaRefAuto{F\colon A\to B,\ x\in A\vdash F(x)\in B}[thm]}
 \proofstep{}{\forall k\in\mathbb N\;
   s[i\mapsto a](k)=\BXLVIIIIf{k=i}{a}{s(k)}}
  {\FormulaRefAuto{TypedTermGraphValues}[thm]{1},
   \FormulaRefAuto{B48Assignments}[def]}
 \proofstep{}{i=i}{\rII}
 \proofstep{}{\BXLVIIIIf{i=i}{a}{s(i)}=a}
  {\FormulaRefAuto{B48IfTrue}[thm]{3}}
 \proofstep{}{s[i\mapsto a](i)=a}{\rIE{4,\rUE{2}}}
 \proofstep{6}{j\neq i}{\rA}
 \proofstep{6}{\BXLVIIIIf{j=i}{a}{s(j)}=s(j)}
  {\FormulaRefAuto{B48IfFalse}[thm]{6}}
 \proofstep{6}{s[i\mapsto a](j)=s(j)}{\rIE{7,\rUE{2}}}
\end{tabproof}

\FormulaThmDeltaK[Uebereinstimmung nach gleichem Belegungswechsel]
{s\equiv_{U\setminus\{i\}}t\ \vdash\
 s[i\mapsto a]\equiv_U t[i\mapsto a]}
{B48AgreementUpdate}{
 \DeltaRow{Belegungen}{s,t\in A_D}
 \DeltaRow{Bereich, Index und Wert}{U\subseteq\mathbb N,\quad i\in\mathbb N,\quad a\in D}
}
\begin{tabproof}
 \proofstep{1}{s\equiv_{U\setminus\{i\}}t}{\rA}
 \proofstep{2}{j\in U}{\rA}
 \proofstep{}{j=i\lor j\neq i}{\FormulaRefAuto{P\lor\neg P}[thm]}
 \proofstep{4}{j=i}{\rA}
 \proofstep{4}{s[i\mapsto a](j)=a}
  {\FormulaRefAuto{B48UpdateValues}[thm]{4}}
 \proofstep{4}{t[i\mapsto a](j)=a}
  {\FormulaRefAuto{B48UpdateValues}[thm]{4}}
 \proofstep{4}{s[i\mapsto a](j)=t[i\mapsto a](j)}
  {\FormulaRefAuto{a=b\dsep c=b\vdash a=c}[thm]{5,6}}
 \proofstep{8}{j\neq i}{\rA}
 \proofstep{2,8}{j\in U\setminus\{i\}}
  {\FormulaRefAuto{B48MemberDifferenceSingleton}[thm]{2,8}}
 \proofstep{1,2,8}{s(j)=t(j)}
  {\FormulaRefAuto{B48Agreement}[def]{1,9}}
 \proofstep{8}{s[i\mapsto a](j)=s(j)}
  {\FormulaRefAuto{B48UpdateValues}[thm]{8}}
 \proofstep{8}{t[i\mapsto a](j)=t(j)}
  {\FormulaRefAuto{B48UpdateValues}[thm]{8}}
 \proofstep{1,2,8}{s[i\mapsto a](j)=t[i\mapsto a](j)}
  {\FormulaRefAuto{a=b\dsep c=b\vdash a=c}[thm]{\rIE{10,11},12}}
 \proofstep{1,2}{s[i\mapsto a](j)=t[i\mapsto a](j)}
  {\rOE{3,4:7,8:13}}
 \proofstep{1}{\forall j\in U\;
                 s[i\mapsto a](j)=t[i\mapsto a](j)}
  {\rUI{\rRI{2,14}}}
 \proofstep{1}{s[i\mapsto a]\equiv_U t[i\mapsto a]}
  {\FormulaRefAuto{B48Agreement}[def]{15}}
\end{tabproof}

\FormulaDefDeltaK[Induktionsbehauptung des Koinzidenzlemmas]
{\begin{aligned}
 C(p)\coloneq\forall s,t\in A_D\;\bigl(
 s\equiv_{\operatorname{FV}(p)}t\ \to{}\\
 ((M,s\models p)\leftrightarrow(M,t\models p))\bigr)
\end{aligned}}
{B48CoincidenceProperty}{
 \DeltaRow{Struktur und Formelcode}{\mathsf{Str}(M),\quad p\in\mathcal F}
}

\FormulaThmDeltaK[Koinzidenzlemma]
{s\equiv_{\operatorname{FV}(p)}t\ \vdash\
 ((M,s\models p)\leftrightarrow(M,t\models p))}
{B48Coincidence}{
 \DeltaRow{Struktur}{\mathsf{Str}(M)}
 \DeltaRow{Formel und Belegungen}{p\in\mathcal F,\quad s,t\in A_D}
}
\begin{tabproofsplitwide}
 \proofpartwideindR[Atomare Formeln]{C(\mathsf e(i,j))\land
   C(\mathsf m(i,j))\land C(\mathsf b)}
  {C(\mathsf e(i,j))\land C(\mathsf m(i,j))\land C(\mathsf b)}
  [B48CoincidenceAtoms]
 \proofstepwidestar[1]{s\equiv_{\{i,j\}}t}{\rA}
 \proofstepwidestar[1]{s(i)=t(i)}
  {\FormulaRefAuto{B48Agreement}[def]{1}}
 \proofstepwidestar[1]{s(j)=t(j)}
  {\FormulaRefAuto{B48Agreement}[def]{1}}
 \proofstepwidestar[]{(s(i)=s(j))\leftrightarrow(s(i)=s(j))}
  {\FormulaRefAuto{P\leftrightarrow P}[thm]}
 \proofstepwidestar[1]{(s(i)=s(j))\leftrightarrow(t(i)=t(j))}
  {\rIE{2,3,4}}
 \proofstepwidestar[1]{\begin{aligned}
 (M,s\models\mathsf e(i,j))\leftrightarrow
 (M,t\models\mathsf e(i,j))
 \end{aligned}}
  {\FormulaRefAuto{B48Satisfaction}[def]{5}}
 \proofstepwidestar[]{((s(i),s(j))\in E)\leftrightarrow
                        ((s(i),s(j))\in E)}
  {\FormulaRefAuto{P\leftrightarrow P}[thm]}
 \proofstepwidestar[1]{((s(i),s(j))\in E)\leftrightarrow
                        ((t(i),t(j))\in E)}
  {\rIE{2,3,7}}
 \proofstepwidestar[1]{\begin{aligned}
 (M,s\models\mathsf m(i,j))\leftrightarrow
 (M,t\models\mathsf m(i,j))
 \end{aligned}}
  {\FormulaRefAuto{B48Satisfaction}[def]{8}}
 \proofstepwidestar[]{C(\mathsf e(i,j))}
  {\FormulaRefAuto{B48CoincidenceProperty}[def]
   {\rUI{\rRI{1,6}}},\FormulaRefAuto{B48FVClauses}[thm]}
 \proofstepwidestar[]{C(\mathsf m(i,j))}
  {\FormulaRefAuto{B48CoincidenceProperty}[def]
   {\rUI{\rRI{1,9}}},\FormulaRefAuto{B48FVClauses}[thm]}
 \proofstepwidestar[]{(M,s\models\mathsf b)\leftrightarrow
                        (M,t\models\mathsf b)}
  {\FormulaRefAuto{B48Satisfaction}[def],
   \FormulaRefAuto{P\leftrightarrow P}[thm]}
 \proofstepwidestar[]{C(\mathsf b)}
  {\FormulaRefAuto{B48CoincidenceProperty}[def]{\rUI{12}}}
 \proofstepwidestar[]{C(\mathsf e(i,j))\land
                      C(\mathsf m(i,j))\land C(\mathsf b)}
  {\rAI{10,\rAI{11,13}}}
 \closeproofpartwideindR

 \proofpartwideindR[Negation]{C(p)\vdash C(\mathsf{neg}(p))}
  {C(p)\vdash C(\mathsf{neg}(p))}[B48CoincidenceNeg]
 \proofstepwidestar[1]{C(p)}{\rA}
 \proofstepwidestar[2]{s\equiv_{\operatorname{FV}(\mathsf{neg}(p))}t}{\rA}
 \proofstepwidestar[2]{s\equiv_{\operatorname{FV}(p)}t}
  {\FormulaRefAuto{B48FVClauses}[thm]{2}}
 \proofstepwidestar[1,2]{(M,s\models p)\leftrightarrow(M,t\models p)}
  {\FormulaRefAuto{B48CoincidenceProperty}[def]{1,3}}
 \proofstepwidestar[1,2]{\neg(M,s\models p)\leftrightarrow
                       \neg(M,t\models p)}
  {\rLRS{4,\FormulaRefAuto{P\leftrightarrow P}[thm]}}
 \proofstepwidestar[1,2]{(M,s\models\neg p)\leftrightarrow
                       (M,t\models\neg p)}
  {\FormulaRefAuto{B48Satisfaction}[def]{5}}
 \proofstepwidestar[1]{C(\mathsf{neg}(p))}
  {\FormulaRefAuto{B48CoincidenceProperty}[def]{\rUI{\rRI{2,6}}}}
 \closeproofpartwideindR

 \proofpartwideindR[Konjunktion]{C(p),C(q)\vdash C(\mathsf{and}(p,q))}
  {C(p),C(q)\vdash C(\mathsf{and}(p,q))}[B48CoincidenceAnd]
 \proofstepwidestar[1]{C(p)}{\rA}
 \proofstepwidestar[2]{C(q)}{\rA}
 \proofstepwidestar[3]{s\equiv_{\operatorname{FV}(\mathsf{and}(p,q))}t}{\rA}
 \proofstepwidestar[3]{s\equiv_{\operatorname{FV}(p)}t}
  {\FormulaRefAuto{B48Agreement}[def]{3},
   \FormulaRefAuto{B48FVClauses}[thm]}
 \proofstepwidestar[3]{s\equiv_{\operatorname{FV}(q)}t}
  {\FormulaRefAuto{B48Agreement}[def]{3},
   \FormulaRefAuto{B48FVClauses}[thm]}
 \proofstepwidestar[1,3]{(M,s\models p)\leftrightarrow(M,t\models p)}
  {\FormulaRefAuto{B48CoincidenceProperty}[def]{1,4}}
 \proofstepwidestar[2,3]{(M,s\models q)\leftrightarrow(M,t\models q)}
  {\FormulaRefAuto{B48CoincidenceProperty}[def]{2,5}}
 \proofstepwidestar[1,2,3]{\begin{aligned}
 ((M,s\models p)\land(M,s\models q))\leftrightarrow{}\\
 ((M,t\models p)\land(M,t\models q))
 \end{aligned}}
  {\rLRS{6,\rLRS{7,\FormulaRefAuto{P\leftrightarrow P}[thm]}}}
 \proofstepwidestar[1,2,3]{(M,s\models p\land q)\leftrightarrow
                          (M,t\models p\land q)}
  {\FormulaRefAuto{B48Satisfaction}[def]{8}}
 \proofstepwidestar[1,2]{C(\mathsf{and}(p,q))}
  {\FormulaRefAuto{B48CoincidenceProperty}[def]{\rUI{\rRI{3,9}}}}
 \closeproofpartwideindR

 \proofpartwideindR[Existenzquantor]{C(p)\vdash C(\mathsf{ex}_i(p))}
  {C(p)\vdash C(\mathsf{ex}_i(p))}[B48CoincidenceExists]
 \proofstepwidestar[1]{C(p)}{\rA}
 \proofstepwidestar[2]{s\equiv_{\operatorname{FV}(\mathsf{ex}_i(p))}t}{\rA}
 \proofstepwidestar[2]{s\equiv_{\operatorname{FV}(p)\setminus\{i\}}t}
  {\FormulaRefAuto{B48FVClauses}[thm]{2}}
 \proofstepwidestar[4]{a\in D}{\rA}
 \proofstepwidestar[4]{s[i\mapsto a]\in A_D}
  {\FormulaRefAuto{B48AssignmentUpdate}[thm]{4}}
 \proofstepwidestar[4]{t[i\mapsto a]\in A_D}
  {\FormulaRefAuto{B48AssignmentUpdate}[thm]{4}}
 \proofstepwidestar[2,4]{s[i\mapsto a]\equiv_{\operatorname{FV}(p)}
                                  t[i\mapsto a]}
  {\FormulaRefAuto{B48AgreementUpdate}[thm]{3,4}}
 \proofstepwidestar[1,2,4]{\begin{aligned}
 (M,s[i\mapsto a]\models p)\leftrightarrow
 (M,t[i\mapsto a]\models p)
 \end{aligned}}
  {\FormulaRefAuto{B48CoincidenceProperty}[def]{1,5,6,7}}
 \proofstepwidestar[1,2]{\begin{aligned}
 \forall a\in D\;((M,s[i\mapsto a]\models p)\leftrightarrow{}\\
 (M,t[i\mapsto a]\models p))
 \end{aligned}}
  {\rUI{\rRI{4,8}}}
 \proofstepwidestar[1,2]{\begin{aligned}
 (\exists a\in D\;(M,s[i\mapsto a]\models p))\leftrightarrow{}\\
 (\exists a\in D\;(M,t[i\mapsto a]\models p))
 \end{aligned}}
  {\FormulaRefAuto{B48ExistsIff}[thm]{9}}
 \proofstepwidestar[1,2]{(M,s\models\exists v_i\,p)\leftrightarrow
                        (M,t\models\exists v_i\,p)}
  {\FormulaRefAuto{B48Satisfaction}[def]{10}}
 \proofstepwidestar[1]{C(\mathsf{ex}_i(p))}
  {\FormulaRefAuto{B48CoincidenceProperty}[def]{\rUI{\rRI{2,11}}}}
 \closeproofpartwideindR

 \proofpartwide{Abschluss der Formelinduktion}
 \proofstepwidestar[]{\{p\in\mathcal F\mid C(p)\}\subseteq\mathcal T}
  {\FormulaRefAuto{B48FormulaInTree}[thm]}
 \proofstepwidestar[]{\forall c\in\Sigma_{\rm at}\;
    \mathsf l(c)\in\{p\in\mathcal F\mid C(p)\}}
  {\FormulaRefAuto{B48CoincidenceAtoms}[thm],
   \FormulaRefAuto{B48FormulaAtom}[thm]}
 \proofstepwidestar[]{\forall p\in\mathcal F\;
     (C(p)\to C(\mathsf{neg}(p)))}
  {\FormulaRefAuto{B48CoincidenceNeg}[thm]}
 \proofstepwidestar[]{\forall p,q\in\mathcal F\;
     (C(p)\land C(q)\to C(\mathsf{and}(p,q)))}
  {\FormulaRefAuto{B48CoincidenceAnd}[thm]}
 \proofstepwidestar[]{\forall i\in\mathbb N\,\forall p\in\mathcal F\;
     (C(p)\to C(\mathsf{ex}_i(p)))}
  {\FormulaRefAuto{B48CoincidenceExists}[thm]}
 \proofstepwidestar[]{\mathsf{Cl}(\{p\in\mathcal F\mid C(p)\})}
  {\FormulaRefAuto{B48FormulaClosure}[def]{1,2,3,4,5},
   \FormulaRefAuto{B48FormulaNeg}[thm],
   \FormulaRefAuto{B48FormulaAnd}[thm],
   \FormulaRefAuto{B48FormulaExists}[thm]}
 \proofstepwidestar[]{\forall p\in\mathcal F\;C(p)}
  {\FormulaRefAuto{B48FormulaInduction}[thm]{6}}
 \proofstepwidestar[8]{s\equiv_{\operatorname{FV}(p)}t}{\rA}
 \proofstepwidestar[8]{(M,s\models p)\leftrightarrow(M,t\models p)}
  {\FormulaRefAuto{B48CoincidenceProperty}[def]{7,8}}
 \closeproofpartwide
\end{tabproofsplitwide}
Die Induktionsbehauptung quantifiziert ueber \emph{alle} Belegungen.
Daher darf sie im Quantorenschritt auf die beiden geaenderten
Belegungen angewendet werden.

\FormulaThmDeltaK[Eine nicht freie Variable darf neu belegt werden]
{i\notin\operatorname{FV}(p)\ \vdash\
 ((M,s\models p)\leftrightarrow(M,s[i\mapsto a]\models p))}
{B48FreshUpdate}{
 \DeltaRow{Struktur und Belegung}{\mathsf{Str}(M),\quad s\in A_D}
 \DeltaRow{Formel, Index und Wert}{p\in\mathcal F,\quad i\in\mathbb N,\quad a\in D}
}
\begin{tabproof}
 \proofstep{1}{i\notin\operatorname{FV}(p)}{\rA}
 \proofstep{2}{j\in\operatorname{FV}(p)}{\rA}
 \proofstep{3}{j=i}{\rA}
 \proofstep{2,3}{i\in\operatorname{FV}(p)}{\rIE{3,2}}
 \proofstep{1,2,3}{\bot}{\rBI{4,1}}
 \proofstep{1,2}{j\neq i}{\rCI{3,5}}
 \proofstep{1,2}{s[i\mapsto a](j)=s(j)}
  {\FormulaRefAuto{B48UpdateValues}[thm]{6}}
 \proofstep{1,2}{s(j)=s[i\mapsto a](j)}
  {\FormulaRefAuto{a=b\vdash b=a}[thm]{7}}
 \proofstep{1}{s\equiv_{\operatorname{FV}(p)}s[i\mapsto a]}
  {\FormulaRefAuto{B48Agreement}[def]{\rUI{\rRI{2,8}}}}
 \proofstep{}{s[i\mapsto a]\in A_D}
  {\FormulaRefAuto{B48AssignmentUpdate}[thm]}
 \proofstep{1}{(M,s\models p)\leftrightarrow
              (M,s[i\mapsto a]\models p)}
  {\FormulaRefAuto{B48Coincidence}[thm]{9,10}}
\end{tabproof}

\FormulaThmDeltaK[Saetze sind unabhaengig von der Belegung]
{\operatorname{FV}(p)=\varnothing\ \vdash\
 ((M,s\models p)\leftrightarrow(M,t\models p))}
{B48SentenceCoincidence}{
 \DeltaRow{Struktur}{\mathsf{Str}(M)}
 \DeltaRow{Formel und Belegungen}{p\in\mathcal F,\quad s,t\in A_D}
}
\begin{tabproof}
 \proofstep{1}{\operatorname{FV}(p)=\varnothing}{\rA}
 \proofstep{}{s\equiv_\varnothing t}{\FormulaRefAuto{B48Agreement}[def]}
 \proofstep{1}{s\equiv_{\operatorname{FV}(p)}t}{\rIE{1,2}}
 \proofstep{1}{(M,s\models p)\leftrightarrow(M,t\models p)}
  {\FormulaRefAuto{B48Coincidence}[thm]{3}}
\end{tabproof}

